\documentclass[11pt,a4paper]{report}

\usepackage[T1]{fontenc}
\usepackage[utf8]{inputenc}
\usepackage{mathpazo}
\usepackage[margin=1.05in]{geometry}
\usepackage{microtype}
\usepackage{amsmath,amssymb,amsthm}
\usepackage{booktabs}
\usepackage{array}
\usepackage{xcolor}
\usepackage{listings}
\usepackage{enumitem}
\usepackage{tikz}
\usetikzlibrary{positioning,arrows.meta,calc,fit,backgrounds}
\usepackage[hidelinks,colorlinks=true,linkcolor=linkblue,citecolor=linkblue,urlcolor=linkblue]{hyperref}

\definecolor{linkblue}{HTML}{1F4E79}
\definecolor{codebg}{HTML}{F6F8FA}
\definecolor{codeframe}{HTML}{D6DBE1}
\definecolor{kw}{HTML}{A3254A}
\definecolor{cmt}{HTML}{5A7D5A}
\definecolor{str}{HTML}{1F6F8B}
\definecolor{accent}{HTML}{1F4E79}
\definecolor{warn}{HTML}{8A5A00}

% ---------- theorem-like environments ----------
\theoremstyle{definition}
\newtheorem{definition}{Definition}[chapter]
\newtheorem{example}[definition]{Example}
\newtheorem{remark}[definition]{Remark}
\theoremstyle{plain}
\newtheorem{theorem}[definition]{Theorem}
\newtheorem{lemma}[definition]{Lemma}
\newtheorem{invariant}[definition]{Invariant}

% ---------- listings ----------
\lstdefinestyle{sqlstyle}{
  language=SQL,
  morekeywords={ASOF,MATCH_CONDITION,LATERAL,OVER,PARTITION,WINDOW,RETURNING,EXPLAIN,
                ANALYZE,BUFFERS,GENERATED,MATERIALIZED,GUC},
  basicstyle=\ttfamily\footnotesize,
  keywordstyle=\color{kw}\bfseries,
  commentstyle=\color{cmt}\itshape,
  stringstyle=\color{str},
  backgroundcolor=\color{codebg},
  frame=single, rulecolor=\color{codeframe},
  framesep=5pt, xleftmargin=6pt, xrightmargin=0pt,
  breaklines=true, showstringspaces=false,
  columns=fullflexible, keepspaces=true,
}
\lstdefinestyle{cstyle}{
  language=C,
  morekeywords={bool,List,Node,Expr,JoinType,PlanState,TupleTableSlot,ExprState,
                AsofJoinState,JoinExpr,Query,Path,Plan,ListCell,Oid,Datum},
  basicstyle=\ttfamily\footnotesize,
  keywordstyle=\color{kw}\bfseries,
  commentstyle=\color{cmt}\itshape,
  stringstyle=\color{str},
  backgroundcolor=\color{codebg},
  frame=single, rulecolor=\color{codeframe},
  framesep=5pt, xleftmargin=6pt,
  breaklines=true, showstringspaces=false,
  columns=fullflexible, keepspaces=true,
}
\lstdefinestyle{algstyle}{
  language={},
  keywords={for,each,in,while,if,else,emit,return,not,and,or,is,EOF,COPY,EMPTY,NULL},
  keywordstyle=\color{kw}\bfseries,
  morecomment=[l]{\#},
  basicstyle=\ttfamily\footnotesize,
  backgroundcolor=\color{codebg},
  frame=single, rulecolor=\color{codeframe},
  framesep=5pt, xleftmargin=6pt,
  breaklines=false, showstringspaces=false,
  columns=fullflexible, keepspaces=true,
  numbers=left, numberstyle=\tiny\color{gray}, numbersep=8pt,
}
\lstdefinestyle{plain}{
  basicstyle=\ttfamily\footnotesize,
  backgroundcolor=\color{codebg},
  frame=single, rulecolor=\color{codeframe},
  framesep=5pt, xleftmargin=6pt,
  breaklines=true, showstringspaces=false,
  columns=fullflexible, keepspaces=true,
}
\lstset{style=sqlstyle}
\emergencystretch=3em
\hbadness=4000

\newcommand{\code}[1]{\texttt{\small #1}}
\newcommand{\file}[1]{\texttt{\footnotesize #1}}
\newcommand{\asof}{{\normalfont\scshape asof}}

% boxed "key idea"
\newenvironment{keyidea}[1][Key idea]{%
  \par\medskip\noindent
  \begin{tikzpicture}
  \node[draw=accent,line width=0.9pt,fill=accent!4,rounded corners=2pt,
        inner sep=9pt,text width=\dimexpr\linewidth-26pt\relax,align=left] (box) \bgroup
  \textbf{\color{accent}#1.}\ \ignorespaces}
  {\egroup; \end{tikzpicture}\par\medskip}

\title{\vspace{-1.5cm}\Huge \textbf{The ASOF Join}\\[0.3cm]
       \Large A Study Guide for Implementing Temporal\\ Nearest-Match Joins in PostgreSQL}
\author{CS631 Project --- Deepjyoti Kakoti (26M0774)\\
        Debasish Bordoloi (26M0776) \quad Shubh Karman Singh (26M0843)}
\date{\today}

\begin{document}
\maketitle

\begin{abstract}
\noindent
An \asof{} join matches each row of a left relation to the \emph{nearest preceding} row of
a right relation, ordered by time and optionally partitioned by a key. It is the natural
operator for aligning two unsynchronised time series --- trades against quotes, readings
against calibration events, facts against slowly-changing dimensions. DuckDB, ClickHouse,
Snowflake, kdb+ and QuestDB all provide it; PostgreSQL does not, and users fall back on a
hand-written \code{LEFT JOIN LATERAL ... ORDER BY ... LIMIT 1}.

This document is the study companion for a two-phase project that adds \asof{} join to
PostgreSQL~17. It develops the operator from its formal definition, explains why the
existing workaround is both awkward and asymptotically worse than it needs to be, walks
through the stages of PostgreSQL's query processor, and then treats each implementation
phase as a technique worth understanding in its own right: Phase~1 as a
\emph{semantics-preserving source-to-source rewrite}, Phase~2 as a \emph{sort-merge
streaming algorithm} with a proof of correctness and a cost model. The emphasis throughout
is on \emph{why} each technique works, not merely what to type.
\end{abstract}

\tableofcontents

% ============================================================
\chapter{The Problem}
% ============================================================

\section{A motivating scenario}

Consider a market-data warehouse with two tables. \emph{Trades} record executions --- what
was actually bought or sold, at what price. \emph{Quotes} record the best bid and ask
prevailing at an instant --- what the market was \emph{offering} at that moment.

\begin{lstlisting}
CREATE TABLE trades (symbol text, ts timestamptz, price numeric, size int);
CREATE TABLE quotes (symbol text, ts timestamptz, bid numeric, ask numeric);
\end{lstlisting}

Both are append-only streams, both are timestamped, and --- crucially --- \emph{their
timestamps do not line up}. This is not sloppiness in the data; it is inherent. Quotes
change when a market maker revises an offer, which may be thousands of times a second or
not at all for minutes. Trades happen when somebody accepts an offer. They are two
independent processes writing to two independent clocks, so the probability that a trade
carries exactly the same timestamp as some quote is essentially zero.

The analytical question is nearly always the same:

\begin{quote}
\emph{For each trade, what was the market showing at the moment it executed?}
\end{quote}

Answering it requires, for each trade, the \textbf{most recent quote for the same symbol at
or before that trade's timestamp}. Three qualifiers, each load-bearing: \emph{most recent}
(not any), \emph{same symbol} (not any symbol), \emph{at or before} (not after).

\subsection{A worked example}

Table~\ref{tab:worked} gives six trades and seven quotes across two symbols. Read it as two
interleaved timelines per symbol.

\begin{table}[h]
\centering
\small
\begin{tabular}{@{}clll r@{\qquad\qquad}clll rr@{}}
\toprule
\multicolumn{5}{c}{\textbf{trades} ($L$)} & \multicolumn{6}{c}{\textbf{quotes} ($R$)}\\
\cmidrule(r){1-5}\cmidrule(l){6-11}
\# & symbol & ts & & price & \# & symbol & ts & & bid & ask\\
\midrule
$\ell_1$ & AAPL & 09:30:02 & & 189.10 & $r_1$ & AAPL & 09:30:00 & & 189.00 & 189.05\\
$\ell_2$ & AAPL & 09:30:07 & & 189.20 & $r_2$ & AAPL & 09:30:05 & & 189.15 & 189.20\\
$\ell_3$ & AAPL & 09:30:08 & & 189.22 & $r_3$ & AAPL & 09:30:09 & & 189.25 & 189.30\\
$\ell_4$ & MSFT & 09:30:04 & & 410.50 & $r_4$ & AAPL & 09:30:12 & & 189.28 & 189.33\\
$\ell_5$ & MSFT & 09:30:11 & & 410.55 & $r_5$ & MSFT & 09:30:01 & & 410.40 & 410.45\\
$\ell_6$ & TSLA & 09:30:06 & & 242.00 & $r_6$ & MSFT & 09:30:10 & & 410.52 & 410.58\\
         &      &          & &        & $r_7$ & MSFT & 09:30:15 & & 410.60 & 410.66\\
\bottomrule
\end{tabular}
\caption{Six trades, seven quotes, three symbols. Note that TSLA has trades but no quotes
at all, and that AAPL quote $r_3$ arrives \emph{between} two AAPL trades.}
\label{tab:worked}
\end{table}

The correct \asof{} left join --- one output row per trade, no more and no fewer --- is
Table~\ref{tab:workedout}.

\begin{table}[h]
\centering
\small
\begin{tabular}{@{}cllrcll p{5.1cm}@{}}
\toprule
\# & symbol & trade ts & price & match & bid & ask & \textbf{why}\\
\midrule
$\ell_1$ & AAPL & 09:30:02 & 189.10 & $r_1$ & 189.00 & 189.05 & only AAPL quote so far\\
$\ell_2$ & AAPL & 09:30:07 & 189.20 & $r_2$ & 189.15 & 189.20 & $r_1$ is stale, $r_3$ is in the future\\
$\ell_3$ & AAPL & 09:30:08 & 189.22 & $r_2$ & 189.15 & 189.20 & \emph{same} match as $\ell_2$: no quote in between\\
$\ell_4$ & MSFT & 09:30:04 & 410.50 & $r_5$ & 410.40 & 410.45 & nearest MSFT quote; AAPL quotes irrelevant\\
$\ell_5$ & MSFT & 09:30:11 & 410.55 & $r_6$ & 410.52 & 410.58 & $r_7$ at 09:30:15 has not happened yet\\
$\ell_6$ & TSLA & 09:30:06 & 242.00 & $\bot$ & NULL & NULL & no TSLA quote exists at all\\
\bottomrule
\end{tabular}
\caption{The answer. Six trades in, six rows out.}
\label{tab:workedout}
\end{table}

\subsection{Reading the answer row by row}

Every row of Table~\ref{tab:workedout} illustrates something the implementation must get
right, so it is worth going through them slowly.

\paragraph{$\ell_1$ --- the straightforward case.} The AAPL trade at 09:30:02 has exactly
one AAPL quote before it, $r_1$ at 09:30:00. Nothing to choose between.

\paragraph{$\ell_2$ --- the case that defines the operator.} The trade at 09:30:07 has two
AAPL quotes before it ($r_1$ at :00, $r_2$ at :05) and two after ($r_3$ at :09, $r_4$ at
:12). The answer is $r_2$. Both wrong answers are instructive:
\begin{itemize}[nosep]
\item Picking $r_1$ gives a \emph{stale} quote --- the market had already moved from
      189.00/189.05 to 189.15/189.20 two seconds before the trade. Any slippage or spread
      analysis built on it is measuring a market state that no longer existed.
\item Picking $r_3$ (09:30:09) means using information from \emph{two seconds after} the
      trade executed.
\end{itemize}

\paragraph{$\ell_3$ --- the repeated match.} The trade at 09:30:08 also matches $r_2$,
because no new AAPL quote arrived between 09:30:07 and 09:30:08. Two distinct left rows
legitimately share one right row. This is normal and common --- quotes are ``state'', and
state persists until it changes --- but it tells us something important about the
implementation: whatever algorithm we use must be able to hand out the same right tuple
more than once, without re-finding it and without consuming it. Chapter~\ref{ch:phase2}
turns exactly this observation into the reason no rewind is needed.

\paragraph{$\ell_4$ --- why the key matters.} The MSFT trade at 09:30:04 must match $r_5$
(MSFT at 09:30:01), even though AAPL's $r_2$ at 09:30:05 is \emph{closer in time} and AAPL's
$r_1$ at 09:30:00 is also nearby. Nearness in time is only meaningful \emph{within} a
symbol; a quote for a different instrument carries no information about this one. The
partition key is not an optimisation, it is part of the definition.

\paragraph{$\ell_5$ --- the boundary.} The MSFT trade at 09:30:11 matches $r_6$ (09:30:10),
one second earlier, and must not match $r_7$ (09:30:15), four seconds later. Note how little
separates a right answer from a wrong one here: a comparison written in the wrong direction
produces $r_7$, which is plausible-looking output with no error and no warning.

\paragraph{$\ell_6$ --- no match at all.} The TSLA trade has no TSLA quote anywhere in the
table. Under \asof{} \emph{left} join it still appears, with NULL bid and ask --- we do not
silently lose a trade just because we cannot price it. Under \asof{} \emph{inner} join it
would be dropped. The choice matters: a trade count that changes depending on quote coverage
is a reconciliation nightmare, which is why the left variant is the common default.

\subsection{Three ways to get this wrong}

The row-by-row walk exposes a hierarchy of failure modes, and they are not equally bad.

\paragraph{1. Staleness (bad).} Matching too old a quote --- $r_1$ instead of $r_2$ for
$\ell_2$. The result is wrong but in a bounded, noisy way: measurements are blurred by
however long the stale period was. It tends to show up as inexplicable variance.

\paragraph{2. Cross-key leakage (worse).} Matching a row from the wrong partition --- AAPL's
quote against an MSFT trade. The result is not merely imprecise, it is meaningless, and
because the timestamps look sensible it can survive a casual review. In the implementation
this failure mode has a specific cause: forgetting to reset per-key state when the key
changes (remark~(B) of \S\ref{sec:candidate}).

\paragraph{3. Lookahead (worst).} Matching a quote from \emph{after} the trade --- $r_3$ for
$\ell_2$, or $r_7$ for $\ell_5$. This is the one failure that is actively dangerous rather
than merely wrong, because it \emph{improves} the apparent quality of the analysis. A
strategy backtested against quotes that were not yet knowable will appear to predict the
market, since it has been handed the answer. The error is invisible in the output --- every
row is populated, every timestamp is nearby, nothing is NULL --- and only shows up as money
lost when the same strategy runs live against quotes that genuinely are not available yet.
In finance this is \emph{lookahead bias}, and avoiding it is precisely why the operator is
defined with $\le$ rather than ``nearest in absolute value''.

\begin{keyidea}[Why not just ``nearest'']
A tempting simplification is ``match the closest quote in either direction''. Apply it to
$\ell_3$, the AAPL trade at 09:30:08: the nearest AAPL quote is $r_3$ at 09:30:09, one second
\emph{after} the trade, beating $r_2$ at 09:30:05 three seconds before it. Symmetric nearness
picks the future quote whenever the next revision happens to be sooner than the last one was
--- which, for any actively quoted instrument, is roughly half the time. \asof{} is
directional by design, and that directionality is the entire point.
\end{keyidea}

\subsection{The same shape outside finance}

Strip away the vocabulary and the pattern is: one table of \textbf{events}, one table of
\textbf{state that changes over time}, joined on ``which state was in force when this event
happened''.

\paragraph{IoT and telemetry.} A sensor reports a reading every second; a technician
recalibrates the sensor occasionally. To interpret a reading you need the calibration in
force at that moment.
\begin{center}\small
\begin{tabular}{@{}llr@{\qquad}llr@{}}
\toprule
\multicolumn{3}{c}{\textbf{readings} (events)} & \multicolumn{3}{c}{\textbf{calibrations} (state)}\\
\cmidrule(r){1-3}\cmidrule(l){4-6}
sensor & ts & raw & sensor & ts & offset\\
\midrule
s7 & 14:02:11 & 512 & s7 & 09:00:00 & $+1.2$\\
s7 & 14:02:12 & 514 & s7 & 13:45:00 & $-0.4$\\
\bottomrule
\end{tabular}
\end{center}
Both readings take the $-0.4$ offset, not $+1.2$: the recalibration at 13:45 superseded it.
Note that the state table is \emph{sparse} --- two rows for a whole day --- while the event
table is dense. That asymmetry ($|L| \gg |R|$ or the reverse) is typical and matters for
costing (Chapter~\ref{ch:planner}).

\paragraph{Slowly-changing dimensions.} A warehouse keeps every version of a customer
record with a validity start date. A sale from March must join the address, tier and
discount the customer had \emph{in March}, not the ones they have today. Using today's
version is the warehouse equivalent of lookahead bias, and \asof{} is precisely the
``as-of-then'' lookup the name refers to.

\paragraph{Experiment and consent analysis.} Each event must be attributed to the most
recent variant assignment, feature-flag value, or consent record preceding it. Attributing
an event to an assignment made \emph{after} it inflates the measured effect of the
experiment --- lookahead bias again, in a different costume.

\paragraph{Operational log correlation.} Pair each error with the last deploy before it. The
deploy after the error did not cause it, however close in time.

In all of these, one side is events, the other is state, and the join condition is
\emph{temporal precedence plus recency}, not equality.

\section{Formal definition}

Fix two relations: $L$ (the \emph{probe} or outer side) and $R$ (the \emph{build} or inner
side). Let $t_L$ and $t_R$ be the match attributes --- the timestamps --- and let
$k_L, k_R$ be (possibly empty) tuples of partition-key attributes.

\begin{definition}[\asof{} match]\label{def:asof}
For a left tuple $\ell \in L$, define the \emph{candidate set}
\[
  C(\ell) \;=\; \{\, r \in R \;:\; k_R(r) = k_L(\ell) \ \wedge\ t_R(r) \le t_L(\ell) \,\}.
\]
The \asof{} match of $\ell$ is
\[
  m(\ell) \;=\;
  \begin{cases}
    \arg\max_{r \in C(\ell)} t_R(r), & C(\ell) \neq \varnothing,\\[2pt]
    \bot, & C(\ell) = \varnothing,
  \end{cases}
\]
where $\bot$ denotes ``no match''. The \emph{\asof{} left join} is
\[
  L \mathbin{\raisebox{0.1ex}{$\bowtie$}_{\!\text{asof}}^{\;\text{left}}} R
  \;=\; \{\, \ell \circ m(\ell) \;:\; \ell \in L \,\},
\]
with $m(\ell)$ replaced by a tuple of NULLs when $m(\ell) = \bot$; the \emph{inner} variant
drops those rows instead.
\end{definition}

\subsection{Reading the definition against the example}

Definition~\ref{def:asof} is compact, so it is worth decoding it against
Table~\ref{tab:worked} one symbol at a time. Take $\ell = \ell_2$, the AAPL trade at
09:30:07.

\begin{itemize}[leftmargin=*]
\item $k_L(\ell_2) = \text{AAPL}$ and $t_L(\ell_2) = \text{09:30:07}$. These are the two
      things about the left row the operator looks at; the price is carried along but plays
      no part in matching.
\item The \textbf{candidate set} $C(\ell_2)$ applies two filters to $R$ at once. The first,
      $k_R(r) = k_L(\ell)$, keeps only AAPL quotes: $\{r_1, r_2, r_3, r_4\}$, discarding the
      MSFT rows. The second, $t_R(r) \le t_L(\ell)$, keeps only those at or before 09:30:07,
      leaving $C(\ell_2) = \{r_1, r_2\}$. Everything the trade \emph{could} legitimately
      match is now in this set, and everything it must not match has been excluded.
\item $\arg\max_{r \in C(\ell)} t_R(r)$ then picks the one with the largest timestamp:
      $r_2$ at 09:30:05. Read ``$\arg\max$'' as ``the element that maximises'' --- it returns
      the \emph{tuple} $r_2$, not the timestamp 09:30:05. This is the ``most recent'' part of
      ``most recent quote at or before''.
\item For $\ell_6$, the TSLA trade, the key filter alone empties the set:
      $C(\ell_6) = \varnothing$, so $m(\ell_6) = \bot$ and the left join emits the trade with
      NULLs.
\item $\ell \circ m(\ell)$ denotes concatenating the left tuple with its match, which is how
      one row of Table~\ref{tab:workedout} is formed.
\end{itemize}

The structure to notice is that the definition is \emph{filter, then take the maximum}, in
that order. Chapter~\ref{ch:phase1} shows that reversing those two steps --- taking the
maximum and then filtering --- gives a different and wrong answer, which is the single most
important correctness point in the whole implementation.

Three properties of Definition~\ref{def:asof} drive everything that follows.

\begin{enumerate}[label=(P\arabic*),leftmargin=*]
\item \textbf{Exactly one match per left row.} The output cardinality of the left variant is
      precisely $|L|$, independent of $|R|$. This is unusual: it makes \asof{} a
      \emph{function} from $L$ to $R \cup \{\bot\}$ rather than a relation, and it means
      the planner can estimate output rows with certainty.
\item \textbf{Asymmetry.} $\arg\max$ picks the \emph{latest preceding} right row. Swapping
      $L$ and $R$ gives a different answer, so \asof{} is \emph{not commutative} --- a fact
      that constrains join reordering (Chapter~\ref{ch:planner}).
\item \textbf{Order dependence.} The operator is defined in terms of $\le$ and $\max$ on
      $t$, so it only makes sense for a type with a total order --- in PostgreSQL terms, a
      type with a \code{btree} operator family. Hash equality is not enough.
\end{enumerate}

\begin{remark}[The four directions]
Definition~\ref{def:asof} uses $\le$ (``the latest quote at or before the trade''), which
is the common case. Three variants follow by symmetry: $<$ (strictly before), $\ge$ and $>$
(the \emph{earliest following} row, i.e.\ $\arg\min$). Our syntax admits all four:
\code{MATCH\_CONDITION (l.ts >= r.ts)} is $\le$ read from the left row's point of view.
Everything in this document is stated for $\le$; the others are obtained by reversing the
comparison and the sort direction, and by excluding the equality case for the strict forms.
\end{remark}

\section{Why ordinary joins cannot express this}

\subsection*{Equi-join is too strict}

\begin{lstlisting}
SELECT * FROM trades t JOIN quotes q
  ON t.symbol = q.symbol AND t.ts = q.ts;   -- almost always empty
\end{lstlisting}
Two independent streams rarely share exact timestamps. With microsecond resolution the
match probability is essentially zero, and the result is an empty relation --- a silent
wrong answer rather than an error.

\subsection*{Range join is too permissive}

Relaxing equality to inequality finds \emph{every} preceding quote, not the nearest one:

\begin{lstlisting}
SELECT * FROM trades t JOIN quotes q
  ON t.symbol = q.symbol AND q.ts <= t.ts;   -- many-to-many blow-up
\end{lstlisting}

The damage is quantitative, not just semantic. Suppose $n$ trades and $m$ quotes are spread
uniformly over the same interval, within a single symbol. A trade at quantile $u \in [0,1]$
of the interval matches roughly $um$ quotes, so the expected intermediate cardinality is
\[
  \sum_{i=1}^{n} \mathbb{E}[u_i m] \;=\; n\,m \int_0^1 u \,du \;=\; \frac{nm}{2}.
\]
For $n = m = 10^6$ that is $5\times 10^{11}$ intermediate rows, to be reduced back to
$10^6$. The work is quadratic in the data size although the answer is linear in it.

\subsection*{Filtering the range join with a window function}

The standard SQL repair picks the nearest row out of the blow-up:

\begin{lstlisting}
SELECT * FROM (
  SELECT t.*, q.bid, q.ask,
         row_number() OVER (PARTITION BY t.symbol, t.ts ORDER BY q.ts DESC) AS rn
  FROM trades t JOIN quotes q
    ON t.symbol = q.symbol AND q.ts <= t.ts
) s WHERE rn = 1;
\end{lstlisting}

This is \emph{correct}, and it is how the operator is often emulated in systems without it.
But it materialises the $nm/2$ intermediate rows before discarding all but $n$ of them: the
filter is applied after the damage. No amount of planner cleverness recovers the linear
algorithm from this formulation, because the window function is opaque to the join planner.

\begin{keyidea}
\asof{} is not sugar. The three SQL encodings above are respectively wrong, quadratic, and
quadratic-then-filtered. A dedicated operator is what allows the engine to exploit
(P1)--(P3) and compute the answer in one ordered pass.
\end{keyidea}


% ============================================================
\chapter{How Other Systems Do It}
% ============================================================

Borrowing a design is cheaper than inventing one, and divergence from established
semantics is a bug report waiting to happen. Table~\ref{tab:systems} summarises the
landscape.

\begin{table}[h]
\centering
\footnotesize
\renewcommand{\arraystretch}{1.25}
\begin{tabular}{@{}p{1.9cm}p{5.1cm}p{2.3cm}p{3.1cm}@{}}
\toprule
\textbf{System} & \textbf{Syntax} & \textbf{Keys} & \textbf{Notes}\\
\midrule
DuckDB &
\code{ASOF JOIN r ON l.k = r.k AND l.ts >= r.ts} &
equalities in \code{ON} &
Inequality lives in \code{ON}; supports \code{LEFT}; one inequality required.\\

ClickHouse &
\code{ASOF LEFT JOIN r ON l.k = r.k AND l.ts >= r.ts} &
equalities in \code{ON} &
Inequality must be last; requires at least one equality key.\\

Snowflake &
\code{ASOF JOIN r MATCH\_CONDITION (l.ts >= r.ts) ON l.k = r.k} &
separate \code{ON} &
Match condition is syntactically distinct from the keys.\\

kdb+/q &
\code{aj[`k`ts; l; r]} &
positional &
Function, not SQL; requires sorted, attribute-tagged tables.\\

QuestDB &
\code{l ASOF JOIN r ON (k)} &
\code{ON (k)} &
Designated timestamp column is implicit in the table definition.\\
\bottomrule
\end{tabular}
\caption{\asof{} join across systems.}
\label{tab:systems}
\end{table}

\section{Why we follow Snowflake}

The project adopts Snowflake's shape:

\begin{lstlisting}
FROM trades t ASOF LEFT JOIN quotes q
  MATCH_CONDITION (t.ts >= q.ts)
  ON t.symbol = q.symbol
\end{lstlisting}

Three reasons, all of them implementation-driven:

\begin{enumerate}[leftmargin=*]
\item \textbf{Separation of concerns.} The DuckDB/ClickHouse style mixes one inequality and
      $n$ equalities in a single \code{ON} clause, then requires the analyser to pick the
      inequality apart from the rest and reject anything else. Keeping
      \code{MATCH\_CONDITION} separate turns a validation problem into a grammar problem,
      and grammar problems produce better error messages and cleaner positions
      (\code{parser\_errposition}).
\item \textbf{No overloading of \code{ON}.} In PostgreSQL, \code{ON} already feeds
      \code{transformJoinOnClause} and ultimately the planner's restriction-clause
      machinery, where an inequality would be treated as an ordinary filter. Giving the
      match condition its own clause means it can be stored in its own \code{JoinExpr}
      field and never mistaken for a qual.
\item \textbf{Room for extension.} A distinct clause gives an obvious home for future
      \code{TOLERANCE}/window bounds (``match only within 5 seconds''), which several
      systems offer and which we may add later.
\end{enumerate}

The cost is that queries are not copy-paste portable from DuckDB. We judge discoverable
errors and a clean parse tree worth more than cross-dialect portability in a teaching
implementation.

\section{Semantic choices to pin down}

Systems quietly disagree on the corners. We must choose, document, and test each:

\begin{table}[h]
\centering
\footnotesize
\renewcommand{\arraystretch}{1.2}
\begin{tabular}{@{}p{4.3cm}p{8.4cm}@{}}
\toprule
\textbf{Corner} & \textbf{Our choice}\\
\midrule
NULL in match column & Never matches (NULL comparisons are unknown, so the row is filtered).\\
NULL in a partition key & Never matches, consistent with \code{=} semantics.\\
Ties on $t_R$ within a key & Implementation-defined: one arbitrary tuple. Documented, and
Phase~2 must agree with Phase~1.\\
Empty right side & Left variant returns all left rows NULL-extended.\\
Strict operators ($<$, $>$) & Equal timestamps excluded.\\
Zero equality keys & Allowed: a single global ordered stream.\\
\bottomrule
\end{tabular}
\caption{Semantic decisions, to be enforced by the regression suite.}
\end{table}

\begin{remark}[On ties]
Picking ``an arbitrary tuple'' is defensible but must be \emph{consistently} arbitrary
across phases, or the differential test of Chapter~\ref{ch:measuring} will flag diffs that
are not bugs. The practical rule: Phase~1's \code{ORDER BY ts DESC LIMIT 1} and Phase~2's
``last candidate seen in ascending order'' happen to agree when the sort is stable on the
same key set, and the tests compare multisets after excluding tie rows.
\end{remark}

% ============================================================
\chapter{What PostgreSQL Does Today}
% ============================================================

\section{The LATERAL workaround}

Without \asof{}, the idiomatic PostgreSQL formulation is a correlated subquery in the
\code{FROM} clause:

\begin{lstlisting}
SELECT t.symbol, t.ts, t.price, q.bid, q.ask
FROM trades t
LEFT JOIN LATERAL (
    SELECT q.bid, q.ask
    FROM quotes q
    WHERE q.symbol = t.symbol
      AND q.ts    <= t.ts
    ORDER BY q.ts DESC
    LIMIT 1
) q ON true;
\end{lstlisting}

\code{LATERAL} is what makes this legal: it permits the subquery to reference \code{t},
the relation to its left. Semantically the subquery is re-evaluated for each row of
\code{trades}, and \code{ORDER BY ... LIMIT 1} performs the $\arg\max$ of
Definition~\ref{def:asof} directly. \code{LEFT JOIN ... ON true} supplies the
NULL-extension when the subquery returns no row.

This is correct. It is also:

\begin{itemize}[nosep]
\item \textbf{verbose} --- eleven lines encoding one idea, with the correlation, the sort
      direction, the \code{LIMIT} and the \code{ON true} each independently easy to get
      wrong;
\item \textbf{fragile} --- writing \code{ORDER BY q.ts} instead of \code{q.ts DESC} yields
      the \emph{oldest} quote, silently;
\item \textbf{dependent on an index} for acceptable performance, as the next section shows.
\end{itemize}

\section{Cost analysis}

The plan is a nested loop with the subquery on the inner side.

\paragraph{With a suitable index.} Given
\code{CREATE INDEX ON quotes (symbol, ts DESC)}, each inner evaluation is a
\code{Limit} over an \code{Index Scan} whose scan key is
\code{symbol = \$1 AND ts <= \$2}. Because the index is ordered, the scan descends to the
right leaf and stops after one tuple. Per outer row the cost is one index descent,
$O(\log m)$, giving
\[
  T_{\text{lateral}} \;=\; O\!\left(n \log m\right).
\]
Each descent is a handful of random page accesses; with $n = 10^6$ that is $10^6$
independent descents, and the constant is dominated by cache misses on the upper index
levels and by per-tuple executor overhead for the \code{Limit}/\code{Index Scan} node pair.

\paragraph{Without one.} The inner side becomes a \code{Sort} over a filtered
\code{Seq Scan} (or at best a \code{Top-N heapsort}), re-executed per outer row:
\[
  T_{\text{lateral}}^{\text{no index}} \;=\; O\!\left(n \cdot m\right),
\]
which for $n = m = 10^6$ is $10^{12}$ tuple comparisons --- hours, not seconds. Nested-loop
rescans are not amortised; PostgreSQL will reuse the inner \code{Materialize} node only
when the inner side is \emph{not} correlated, and here it necessarily is.

\begin{keyidea}
The workaround's shape forces a \emph{per-row search}. Even in the good case it pays
$\log m$ random descents $n$ times, because nothing tells the engine that consecutive
outer rows probe nearby positions in the same ordered inner relation. Phase~2's whole
contribution is to exploit exactly that locality.
\end{keyidea}

\section{A third problem: planning}

Beyond asymptotics, the LATERAL form limits the planner. The subquery is a black box with
a \code{LIMIT}: it cannot be hash-joined, cannot be parallelised across the inner side, and
its selectivity is estimated through the generic subquery path. The optimiser cannot, for
instance, decide to sort both inputs once --- that plan is not in the search space, because
the query as written does not describe it.


% ============================================================
\chapter{Anatomy of a PostgreSQL Query}\label{ch:anatomy}
% ============================================================

To extend the engine one must know where to cut. PostgreSQL processes a query in five
stages, each consuming one data structure and producing the next.

\begin{figure}[h]
\centering
\begin{tikzpicture}[
  node distance=5mm,
  stage/.style={draw=accent,fill=accent!6,rounded corners=2pt,minimum height=9mm,
                inner xsep=5pt,font=\footnotesize\bfseries,align=center},
  arr/.style={-{Stealth[length=2mm]},accent,thick},
  lbl/.style={font=\scriptsize\ttfamily,text=black!65}
]
\node[stage] (p) {Parse};
\node[stage,right=of p] (a) {Analyze};
\node[stage,right=of a] (r) {Rewrite};
\node[stage,right=of r] (o) {Plan};
\node[stage,right=of o] (e) {Execute};
\draw[arr] (p)--(a); \draw[arr] (a)--(r); \draw[arr] (r)--(o); \draw[arr] (o)--(e);
\node[lbl,below=1.5mm of p] {RawStmt};
\node[lbl,below=1.5mm of a] {Query};
\node[lbl,below=1.5mm of r] {Query[]};
\node[lbl,below=1.5mm of o] {PlannedStmt};
\node[lbl,below=1.5mm of e] {tuples};
\node[font=\scriptsize,above=1mm of p,text=black!60] {gram.y};
\node[font=\scriptsize,above=1mm of a,text=black!60] {parse\_clause.c};
\node[font=\scriptsize,above=1mm of r,text=black!60] {rewriteHandler.c};
\node[font=\scriptsize,above=1mm of o,text=black!60] {joinpath.c};
\node[font=\scriptsize,above=1mm of e,text=black!60] {nodeXxx.c};
\end{tikzpicture}
\caption{The five stages, the structure each produces, and a representative source file.}
\label{fig:pipeline}
\end{figure}

\section{Stage by stage}

\paragraph{1. Parse (\file{src/backend/parser/gram.y}, \file{scan.l}).}
A Bison grammar turns text into a \emph{raw parse tree}: a tree of nodes that mirrors the
syntax and resolves \emph{nothing}. Table names are strings, column references are
\code{ColumnRef} lists, and no catalog is consulted. The parser's only job is to decide
whether the text is well-formed and to record its shape. A \code{FROM} clause join becomes
a \code{JoinExpr} with \code{larg}, \code{rarg}, \code{jointype}, \code{quals} and
\code{usingClause}.

\paragraph{2. Analyze (\file{parser/analyze.c}, \file{parse\_clause.c}, \file{parse\_expr.c}).}
Semantic analysis resolves names against the catalog, assigns types, and produces a
\code{Query}. Range-table entries (\code{RangeTblEntry}) replace table names; \code{Var}
nodes with \code{(varno, varattno)} replace column references; operators are resolved to
specific \code{Oid}s through overload resolution. The \code{FROM} clause is processed by
\code{transformFromClauseItem()}, which is recursive over \code{JoinExpr} --- this is the
function Phase~1 hooks.

\paragraph{3. Rewrite (\file{rewrite/rewriteHandler.c}).}
Views are expanded and rules applied, \code{Query} in, \code{Query} out. The \asof{}
project does not touch this stage, but it matters for one reason: a view defined over an
\asof{} join must survive expansion, which is why the node must deparse correctly
(\file{ruleutils.c}).

\paragraph{4. Plan (\file{optimizer/}).}
The planner converts a \code{Query} into a \code{PlannedStmt}. Internally it works with
\code{RelOptInfo} (a relation, base or join) and \code{Path} (one way of producing that
relation, with a cost and a sort order). \code{add\_paths\_to\_joinrel()} in
\file{path/joinpath.c} proposes every join strategy for a pair of relations; the cheapest
surviving \code{Path} is turned into a \code{Plan} by \file{plan/createplan.c}. Phase~2
adds a \code{Path} type here and teaches \file{path/costsize.c} what it costs.

\paragraph{5. Execute (\file{executor/}).}
The \code{Plan} tree is instantiated as a \code{PlanState} tree and pulled one tuple at a
time: each node's \code{ExecProcNode} returns a \code{TupleTableSlot} or NULL at end of
stream. This is the Volcano iterator model, and it is why a streaming merge algorithm fits
so naturally. Phase~2 adds \file{executor/nodeAsofjoin.c}.

\begin{table}[h]
\centering
\footnotesize
\renewcommand{\arraystretch}{1.2}
\begin{tabular}{@{}llll@{}}
\toprule
\textbf{Stage} & \textbf{Structure} & \textbf{Phase 1 touches} & \textbf{Phase 2 touches}\\
\midrule
Parse    & \code{RawStmt}, \code{JoinExpr} & grammar, new fields & --- \\
Analyze  & \code{Query}, \code{Var}, \code{RTE} & validation + rewrite & emit \code{JOIN\_ASOF}\\
Rewrite  & \code{Query} & deparse only & deparse only\\
Plan     & \code{Path}, \code{Plan}, \code{PathKey} & --- & path, cost, plan\\
Execute  & \code{PlanState}, slots & --- & new node\\
\bottomrule
\end{tabular}
\caption{Where each phase intervenes. Phase~1 deliberately stops at stage~2.}
\label{tab:touchpoints}
\end{table}

\begin{keyidea}
Phase~1 and Phase~2 differ in \emph{which stage they stop at}. Phase~1 rewrites the parse
tree and lets stages 3--5 run untouched; Phase~2 carries a new concept all the way to
stage~5. That is why Phase~1 is low-risk and Phase~2 is where the performance is.
\end{keyidea}

\section{Two structures worth studying before Phase 2}

\paragraph{\code{PathKey}.} A \code{PathKey} is the planner's representation of ``this path
produces tuples ordered by expression $e$, using operator family $f$, ascending or
descending, nulls first or last''. A list of them is a sort order. Crucially,
\code{PathKey}s are \emph{canonicalised}: two paths claiming the same order share pointer-equal
\code{PathKey} objects, so order compatibility is tested by list-prefix comparison
rather than expression matching. This is the mechanism by which a merge join avoids a sort
when an index already supplies the order --- and the mechanism Phase~2 reuses.

\paragraph{\code{TupleTableSlot}.} A slot is a holder for one tuple, possibly virtual
(array of \code{Datum}s), possibly a pointer into a buffer page. Slots are
\emph{overwritten} by the next \code{ExecProcNode} call on the same child. Any tuple a node
needs to remember across calls must be copied with \code{ExecCopySlot} into a slot the node
owns --- a rule that Phase~2's candidate buffer depends on absolutely
(\S\ref{sec:candidate}).


% ============================================================
\chapter{Phase 1 --- The Rewrite Technique}\label{ch:phase1}
% ============================================================

\section{Why rewrite at all?}

It is tempting to go straight to the native operator. Rewriting first buys four things:

\begin{enumerate}[leftmargin=*]
\item \textbf{Zero executor risk.} The rewrite produces a plan the engine already knows how
      to run correctly. No new crash surface, no new memory-management bugs.
\item \textbf{A reference implementation.} Phase~2 must agree with \emph{something}. The
      rewrite is an executable specification of Definition~\ref{def:asof}, and
      Chapter~\ref{ch:measuring}'s differential test compares the two implementations
      directly --- an oracle we would otherwise lack.
\item \textbf{Early feedback on syntax.} The grammar, the error messages and the
      documentation are all exercised before any planner work begins. If the syntax is
      wrong, we learn it cheaply.
\item \textbf{A fallback.} Behind a GUC, the rewrite remains available when the native path
      is disabled or inapplicable.
\end{enumerate}

This is a standard compiler technique --- \emph{desugaring}, or source-to-source
translation into a core language. The new construct is defined by the translation, and the
existing machinery does the work.

\section{Grammar mechanics}

\subsection{Keyword categories}

PostgreSQL classifies every keyword into one of four categories in
\file{src/include/parser/kwlist.h}:

\begin{table}[h]
\centering
\footnotesize
\begin{tabular}{@{}lll@{}}
\toprule
\textbf{Category} & \textbf{Usable as identifier?} & \textbf{Example}\\
\midrule
\code{UNRESERVED\_KEYWORD}   & yes                       & \code{ASOF} (ours)\\
\code{COL\_NAME\_KEYWORD}    & as a column name only     & \code{BETWEEN}\\
\code{TYPE\_FUNC\_NAME\_KEYWORD} & not as type/function name & \code{LEFT}\\
\code{RESERVED\_KEYWORD}     & no                        & \code{SELECT}\\
\bottomrule
\end{tabular}
\caption{Keyword categories, from least to most intrusive.}
\end{table}

We register \code{ASOF} and \code{MATCH\_CONDITION} as \textbf{unreserved} and add them to
the \code{unreserved\_keyword:} production in \file{gram.y}. The reason is compatibility:
an existing application may well have a column called \code{asof} or
\code{match\_condition}, and reserving the word would break it. Unreserved keywords are
distinguishable by the parser from context alone.

\begin{keyidea}
Adding a keyword is a compatibility decision, not a syntax decision. Always start at
unreserved and escalate only if the grammar genuinely cannot be made unambiguous ---
and treat ``cannot'' as a claim to be proved by Bison's conflict report, not assumed.
\end{keyidea}

\subsection{The production}

The join syntax lives in the \code{joined\_table} production. The addition is, in essence:

\begin{lstlisting}[style=plain]
joined_table:
    ...
  | table_ref ASOF opt_outer_join_type JOIN table_ref
        MATCH_CONDITION '(' a_expr ')' opt_asof_on
      {
          JoinExpr *n = makeNode(JoinExpr);
          n->jointype       = $3;         /* JOIN_INNER or JOIN_LEFT */
          n->isAsof         = true;
          n->larg           = $1;
          n->rarg           = $5;
          n->matchCondition = $8;
          n->quals          = $10;        /* may be NULL */
          $$ = (Node *) n;
      }
\end{lstlisting}

Two design points deserve comment.

\paragraph{\code{MATCH\_CONDITION} is mandatory in the grammar.} Making it optional and
checking later would produce a worse error (``\asof{} join requires a match condition'',
with no position) than letting Bison report a syntax error at the exact token. Push
validation into the grammar wherever the grammar can express it.

\paragraph{\code{RIGHT}/\code{FULL} are accepted by the grammar and rejected by the
analyser.} This looks backwards, but it produces the better message. If the grammar
excluded them, \code{ASOF RIGHT JOIN} would yield ``syntax error at or near RIGHT'', which
tells the user nothing. Accepting it and then raising \emph{\asof{} join does not support
RIGHT outer joins} with a hint is strictly more helpful. The rule of thumb inverts the
previous one: push validation into the analyser when the analyser can say something the
grammar cannot.

\subsection{Conflicts}

\file{gram.y} is an LALR(1) grammar with a known, non-zero baseline conflict count
(declared via \code{\%expect}). Any \emph{new} shift/reduce or reduce/reduce conflict means
the grammar has become ambiguous at one token of lookahead, and must be fixed rather than
suppressed. Because \code{ASOF} appears immediately after a complete \code{table\_ref} ---
the same position as \code{LEFT}, \code{INNER}, \code{CROSS} --- and is followed by a fixed
token sequence, this addition is conflict-free. Verify, do not assume: build
\file{src/backend/parser} and compare Bison's report against the baseline.

\section{Parse-tree surgery}

\subsection{The translation}

The heart of Phase~1 is a structural rewrite performed in a new
\file{src/backend/parser/parse\_asof.c}, invoked from \code{transformFromClauseItem()}
\emph{before} the right-hand side is transformed. Given

\begin{lstlisting}
L ASOF LEFT JOIN R MATCH_CONDITION (L.ts >= R.ts) ON L.k = R.k
\end{lstlisting}

we construct the parse tree that the following text would have produced:

\begin{lstlisting}
L LEFT JOIN LATERAL (
    SELECT * FROM R
    WHERE R.ts <= L.ts AND R.k = L.k
    ORDER BY R.ts DESC
    LIMIT 1
) ON true
\end{lstlisting}

Concretely: a \code{RangeSubselect} wrapping a \code{SelectStmt} whose \code{fromClause} is
the original \code{rarg}, with \code{lateral = true}, substituted as the \code{rarg} of an
ordinary \code{JoinExpr} with \code{quals = makeBoolConst(true, false)}.

\begin{keyidea}
Rewriting the \emph{raw} parse tree, not the analysed \code{Query}, is what makes this
cheap. Name resolution, type resolution, \code{LATERAL} validity checking, lateral
reference marking and subquery planning are all performed afterwards by existing code that
has no idea \asof{} exists.
\end{keyidea}

\subsection{Direction, strictness, and the four operators}

The match condition is written from the left row's point of view; the subquery's
\code{WHERE} needs it from the right row's. Table~\ref{tab:directions} is the whole of it.

\begin{table}[h]
\centering
\footnotesize
\renewcommand{\arraystretch}{1.25}
\begin{tabular}{@{}llll@{}}
\toprule
\textbf{\code{MATCH\_CONDITION}} & \textbf{Meaning} & \textbf{Subquery \code{WHERE}} & \textbf{\code{ORDER BY}}\\
\midrule
\code{l.ts >= r.ts} & latest at or before & \code{r.ts <= l.ts} & \code{r.ts DESC}\\
\code{l.ts > r.ts}  & latest strictly before & \code{r.ts < l.ts} & \code{r.ts DESC}\\
\code{l.ts <= r.ts} & earliest at or after & \code{r.ts >= l.ts} & \code{r.ts ASC}\\
\code{l.ts < r.ts}  & earliest strictly after & \code{r.ts > l.ts} & \code{r.ts ASC}\\
\bottomrule
\end{tabular}
\caption{The four admissible match conditions and their translations.}
\label{tab:directions}
\end{table}

\subsection{Why the equality keys go \emph{inside} the subquery}

This is the subtlest point in Phase~1, and a natural place to get it wrong. One might
place \code{L.k = R.k} in the outer join's \code{ON} clause instead:

\begin{lstlisting}
L LEFT JOIN LATERAL (
    SELECT * FROM R WHERE R.ts <= L.ts ORDER BY R.ts DESC LIMIT 1
) q ON L.k = q.k            -- WRONG
\end{lstlisting}

The result is wrong, not merely slower. The subquery now computes ``the latest preceding
row of $R$ \emph{over all keys}'', and the outer \code{ON} then discards it if its key
happens to differ --- producing NULL where a correct match exists. The $\arg\max$ of
Definition~\ref{def:asof} is taken over the candidate set $C(\ell)$, which is
\emph{already restricted by the key}; restricting afterwards does not commute with
$\arg\max$.

\begin{lemma}[Key placement]
Let $\sigma_\theta$ be selection and $\mathrm{top}_1^{t}$ the ``highest $t$'' operator.
Then $\mathrm{top}_1^{t}(\sigma_\theta(R)) \neq \sigma_\theta(\mathrm{top}_1^{t}(R))$ in
general: the left side returns the best $\theta$-satisfying row, the right side returns the
globally best row if it satisfies $\theta$ and nothing otherwise.
\end{lemma}

Concretely, take $\ell_4$ (MSFT, 09:30:04) from Table~\ref{tab:worked}. Without the key
predicate, the subquery computes ``the latest quote of \emph{any} symbol at or before
09:30:04''. The candidates are $r_1$ (AAPL, 09:30:00) and $r_5$ (MSFT, 09:30:01); the latest
is $r_5$, the outer \code{ON} accepts it because the symbols match, and the query happens to
produce the right answer.

That accident is what makes the bug dangerous. Add a single AAPL quote at 09:30:03 --- an
entirely ordinary thing for the data to contain --- and the subquery now returns \emph{that}
row as the latest. The outer \code{ON} rejects it, AAPL~$\neq$~MSFT, and $\ell_4$ reports no
match at all: a trade with a perfectly good MSFT quote three seconds before it comes back
NULL. The failure surfaces only when another key's rows interleave, so it passes a small
hand-made test and fails on real data --- precisely the shape of bug the differential test of
Chapter~\ref{ch:measuring} exists to catch.

This lemma is why the predicate must be pushed \emph{through} the \code{LIMIT}, i.e.\ into
the subquery, where it filters before the $\arg\max$. It also explains why PostgreSQL's own
optimiser will never perform this pushdown for us: predicate pushdown past a \code{LIMIT}
is unsound, and \code{allpaths.c} correctly refuses it.

\section{Correctness of the translation}

\begin{theorem}[Rewrite equivalence]\label{thm:rewrite}
For a match condition $t_L \ge t_R$ and key equality $k_L = k_R$, the LATERAL translation
of \S5.3.1 produces exactly the relation of Definition~\ref{def:asof}, up to the choice
among tied $t_R$ values.
\end{theorem}

\begin{proof}[Proof sketch]
Fix $\ell \in L$. By the semantics of \code{LATERAL}, the subquery is evaluated with
$\ell$'s columns bound as constants. Its \code{WHERE} clause then selects exactly
$\{\, r \in R : k_R(r) = k_L(\ell) \wedge t_R(r) \le t_L(\ell) \,\} = C(\ell)$, which is
Definition~\ref{def:asof}'s candidate set. \code{ORDER BY t DESC LIMIT 1} returns a tuple
of $C(\ell)$ with maximal $t_R$ when $C(\ell) \ne \varnothing$ --- i.e.\ an element of
$\arg\max$ --- and the empty relation otherwise. \code{LEFT JOIN ... ON true} pairs $\ell$
with that tuple, or NULL-extends it when the subquery is empty, matching $m(\ell) = \bot$.
Since the construction is per-$\ell$ and the left join preserves every $\ell$ exactly once,
the two relations agree tuple-for-tuple. Where several $r$ share the maximal $t_R$, both
formulations select one arbitrarily; agreement holds up to that choice.
\end{proof}

\begin{remark}
The restriction to \emph{one} top-level comparison in the match condition (enforced by the
analyser) is what makes the proof go through: a conjunction or disjunction of inequalities
would not reduce to a single \code{ORDER BY} direction. This is also exactly the restriction
Phase~2 needs, which is why it is checked in Phase~1 rather than deferred.
\end{remark}

\section{Validation and error messages}

The analyser must verify, with a \code{parser\_errposition()} on the offending node:

\begin{enumerate}[nosep,leftmargin=*]
\item the match condition is a single \code{OpExpr} whose operator is one of
      \code{>=, >, <=, <};
\item its left operand references only $L$ and its right operand only $R$ --- checked with
      \code{pull\_varnos} against the respective namespace items --- or vice versa, in which
      case the comparison is flipped;
\item the operand types share a \code{btree} operator family, so the comparison has a total
      order. This matters for Phase~1 only insofar as \code{ORDER BY} needs it, but it is
      \emph{essential} for Phase~2's merge, so it is rejected early rather than becoming a
      planner-time failure;
\item the \code{ON} clause, if present, is a conjunction of plain $L$-to-$R$ equalities;
\item the join type is \code{INNER} or \code{LEFT}.
\end{enumerate}

Every one of these gets a regression test with stable expected output. An error message is
an interface: once released it is as hard to change as the syntax itself.


% ============================================================
\chapter{Phase 2 --- The Merge Technique}\label{ch:phase2}
% ============================================================

\section{The insight}

\begin{keyidea}
If both inputs are sorted by $(k, t)$ ascending, then as the outer cursor advances, the
\asof{} match also advances monotonically. The nearest preceding inner row is simply
\emph{the last inner row seen before the boundary} --- so a single synchronised pass over
both inputs suffices, with no search and no backtracking.
\end{keyidea}

The LATERAL plan re-searches the inner relation from scratch for every outer row, throwing
away everything it learned from the previous one. But the outer rows arrive in increasing
$t$ order, so each successive probe lands at or after the previous one. Maintaining a
single inner cursor plus one remembered tuple replaces $n$ independent $O(\log m)$
descents with $n + m$ sequential steps.

\section{The algorithm}
\label{sec:candidate}

\begin{lstlisting}[style=algstyle,caption={\asof{} left join, both inputs sorted ascending
by $(k,t)$. \code{cand} is a slot owned by this node.},captionpos=b]
AsofJoin(outer, inner):
  cand        <- EMPTY          # best candidate for the current outer row
  cand_key    <- NULL           # the key group cand belongs to
  r           <- next(inner)

  for each l in outer:                                  # (A) ascending by (k,t)
      if key(l) != cand_key:                            # (B) new key group
          cand     <- EMPTY
          cand_key <- key(l)

      while r is not EOF and key(r) <  key(l):          # (C) skip stale keys
          r <- next(inner)

      while r is not EOF and key(r) == key(l)
                        and t(r) <= t(l):               # (D) absorb candidates
          cand <- COPY(r)                               # (E) must copy!
          r    <- next(inner)

      if cand is EMPTY:
          emit l o NULLs                                # left-outer
      else:
          emit l o cand
\end{lstlisting}

Five remarks, keyed to the comments:

\begin{description}[leftmargin=*,font=\normalfont\bfseries]
\item[(A) Both sides are sorted.] The planner guarantees this by requiring pathkeys
      (Chapter~\ref{ch:planner}); the executor may assume it.
\item[(B) Key groups are contiguous.] Because the sort is on $(k, t)$ with $k$ leading, all
      rows of a key arrive together. A key change resets the candidate --- forgetting this
      reset is the classic bug, and it leaks matches from one symbol to the next.
\item[(C) and (D) The inner cursor never rewinds.] Line~(C) discards inner keys that the
      outer has passed; line~(D) advances while the inner row still satisfies the match
      condition. Both loops only ever call \code{next(inner)}, which is why the total inner
      work is $O(m)$ across the entire join rather than per outer row.
\item[(E) The candidate must be \emph{copied}.] \code{cand <- COPY(r)} is
      \code{ExecCopySlot}, not a pointer assignment. The inner child reuses its output slot
      on the next \code{ExecProcNode} call, so a retained pointer becomes a dangling
      reference to a tuple that has already been overwritten --- a bug that manifests as
      sporadically wrong values rather than a crash, and therefore one to prevent by
      construction.
\item[Strict operators] replace \code{t(r) <= t(l)} with \code{t(r) < t(l)} in~(D).
      The ``earliest following'' variants sort descending, or equivalently reverse the
      comparison --- the loop structure is unchanged.
\end{description}

\subsection{Why no mark/restore}

A merge \emph{equi}-join must handle duplicate join keys on the inner side: when the outer
advances to another row with the same key, the inner cursor has to be rewound to the start
of that key's run. PostgreSQL implements this with \code{ExecMarkPos}/\code{ExecRestrPos}
on the inner node, which in turn forces a \code{Materialize} node when the inner plan cannot
rewind cheaply --- a significant complication in \file{nodeMergejoin.c}.

\asof{} join needs none of it. By property (P1) each outer row has exactly one match, and
that match is \emph{the last candidate absorbed} --- a single tuple, already held in
\code{cand}. Consecutive outer rows with the same key continue absorbing from where the
previous one stopped; the candidate is only ever extended forward, never re-scanned. The
one-tuple buffer replaces the entire mark/restore mechanism.

\begin{keyidea}
\asof{}'s ``exactly one match'' property is not just a semantic curiosity --- it is what
collapses merge join's rewind machinery into a single remembered tuple. A weaker operator
would need the full apparatus.
\end{keyidea}

\section{Correctness}

\begin{invariant}[Candidate]\label{inv:cand}
At the point where the algorithm emits a row for outer tuple $\ell$, either
\code{cand} holds a tuple $r^\star$ with
\[
  r^\star \in C(\ell) \quad\text{and}\quad t_R(r^\star) = \max_{r \in C(\ell)} t_R(r),
\]
or \code{cand} is empty and $C(\ell) = \varnothing$.
\end{invariant}

\begin{proof}[Proof]
By induction over the outer stream.

\emph{Base.} Before the first outer tuple, \code{cand} is empty and \code{cand\_key} is
NULL, so line~(B) resets and the loops at (C),(D) run from the start of the inner stream.

\emph{Step.} Assume the invariant held for the previous outer tuple $\ell'$. Two cases.

\textsc{Case 1: } $k(\ell) \neq k(\ell')$. Line~(B) empties \code{cand}. Line~(C) advances
$r$ past all inner keys less than $k(\ell)$; since the inner stream is sorted, no inner row
with key $k(\ell)$ has been consumed and discarded (consumption only happens in~(C), for
strictly smaller keys, or in~(D), for the previous key group). Line~(D) then absorbs
precisely those rows with key $k(\ell)$ and $t_R \le t_L(\ell)$, in ascending $t_R$ order,
assigning each to \code{cand}; the last assignment is therefore the maximum. If no such row
exists, \code{cand} remains empty and indeed $C(\ell) = \varnothing$.

\textsc{Case 2: } $k(\ell) = k(\ell')$. Then \code{cand} still holds the maximal element of
$C(\ell')$. Because the outer stream is sorted on $(k,t)$, $t_L(\ell) \ge t_L(\ell')$, hence
$C(\ell') \subseteq C(\ell)$: the candidate set only grows. The rows of
$C(\ell) \setminus C(\ell')$ are exactly those with key $k(\ell)$ and
$t_L(\ell') < t_R \le t_L(\ell)$, which are the next rows in the inner stream and are
precisely what loop~(D) now absorbs, again in ascending order. After the loop \code{cand}
holds the maximum over $C(\ell') \cup (C(\ell)\setminus C(\ell')) = C(\ell)$, since the
previous candidate was maximal over $C(\ell')$ and each absorbed row is compared by
recency. If both are empty, $C(\ell) = \varnothing$.
\end{proof}

\begin{theorem}[Algorithm correctness]
Given inputs sorted ascending on $(k,t)$, the algorithm of \S\ref{sec:candidate} computes
the \asof{} left join of Definition~\ref{def:asof}.
\end{theorem}
\begin{proof}
The outer loop visits each $\ell \in L$ exactly once and emits exactly one row for it, so
the output has $|L|$ rows as required by (P1). By Invariant~\ref{inv:cand}, the emitted row
is $\ell \circ m(\ell)$ when $C(\ell) \neq \varnothing$ and $\ell$ NULL-extended otherwise,
which is Definition~\ref{def:asof}. The inner variant differs only in suppressing the
NULL-extended rows.
\end{proof}

\begin{remark}[The monotonicity requirement]
Case~2 of the induction uses $t_L(\ell) \ge t_L(\ell')$ --- the outer stream must be sorted
\emph{within} a key. If it were not, $C$ could shrink between consecutive outer rows and the
retained candidate would be too recent. This is precisely why the planner must enforce a
pathkey on the outer side too, not only on the inner (\S\ref{sec:pathkeys}); sorting only
the inner side gives a plausible-looking plan that is quietly wrong.
\end{remark}

\section{Complexity}

Let $n = |L|$, $m = |R|$.

\[
  T_{\text{asof}} \;=\;
  \underbrace{O(n \log n) + O(m \log m)}_{\text{sorts, if needed}}
  \;+\;
  \underbrace{O(n + m)}_{\text{merge pass}}
\]

and $O(n + m)$ alone when indexes already supply both orders. Compare with
$O(n \log m)$ random descents for the LATERAL plan. The asymptotic gap is modest; the
\emph{constant} is where the win lives:

\begin{itemize}[nosep]
\item the merge pass is sequential, so it is prefetch- and cache-friendly, whereas index
      descents are random;
\item it performs one executor node invocation per tuple, not one subplan execution per
      outer row;
\item sorts are performed once, in bulk, with \code{work\_mem}-sized runs, and are
      skipped entirely when an index on $(k, t)$ exists.
\end{itemize}

Memory is $O(1)$ beyond the sorts: exactly one candidate tuple, regardless of data size ---
in contrast to a hash-based strategy, which is not even applicable here because the
predicate is an inequality.

\section{A worked example}

Figure~\ref{fig:merge} traces the algorithm on one key group.

\begin{figure}[h]
\centering
\begin{tikzpicture}[
  x=1.08cm, y=1cm,
  oin/.style={circle,draw=accent,fill=accent!12,inner sep=1.6pt,font=\scriptsize},
  rin/.style={rectangle,draw=warn,fill=warn!12,inner sep=2.2pt,font=\scriptsize},
  ax/.style={-{Stealth[length=2mm]},black!55},
  mt/.style={-{Stealth[length=1.6mm]},accent!80,dashed,thin}
]
% axes
\draw[ax] (-0.4,2) -- (9.3,2) node[right,font=\scriptsize] {$t$};
\draw[ax] (-0.4,0) -- (9.3,0) node[right,font=\scriptsize] {$t$};
\node[font=\scriptsize\bfseries,left=2mm] at (-0.4,2) {outer $L$};
\node[font=\scriptsize\bfseries,left=2mm] at (-0.4,0) {inner $R$};

% outer tuples at t = 2, 7, 11, 12
\node[oin] (l1) at (1,2) {}; \node[font=\scriptsize,above=0.5mm] at (1,2) {2};
\node[oin] (l2) at (3.5,2) {}; \node[font=\scriptsize,above=0.5mm] at (3.5,2) {7};
\node[oin] (l3) at (6,2) {}; \node[font=\scriptsize,above=0.5mm] at (6,2) {11};
\node[oin] (l4) at (7,2) {}; \node[font=\scriptsize,above=0.5mm] at (7,2) {12};

% inner tuples at t = 1, 5, 6, 10, 14
\node[rin] (r1) at (0.5,0) {}; \node[font=\scriptsize,below=0.5mm] at (0.5,0) {1};
\node[rin] (r2) at (2.5,0) {}; \node[font=\scriptsize,below=0.5mm] at (2.5,0) {5};
\node[rin] (r3) at (3,0) {};   \node[font=\scriptsize,below=0.5mm] at (3,0) {6};
\node[rin] (r4) at (5,0) {};   \node[font=\scriptsize,below=0.5mm] at (5,0) {10};
\node[rin] (r5) at (8,0) {};   \node[font=\scriptsize,below=0.5mm] at (8,0) {14};

% match arrows
\draw[mt] (l1) -- (r1);
\draw[mt] (l2) -- (r3);
\draw[mt] (l3) -- (r4);
\draw[mt] (l4) -- (r4);

\node[font=\scriptsize,text=accent!85,align=center] at (3.6,-1.15)
      {each outer row takes the last inner row to its left};
\node[font=\scriptsize,text=warn,align=center] at (8.3,1.05) {never\\reached};
\end{tikzpicture}
\caption{One key group. Outer timestamps $2,7,11,12$; inner $1,5,6,10,14$. The inner
cursor moves left to right exactly once. Note that $t=5$ is absorbed and then superseded by
$t=6$ --- it is a candidate briefly --- and that outer rows $11$ and $12$ share the
candidate $10$ without any rewind.}
\label{fig:merge}
\end{figure}

\begin{table}[h]
\centering
\footnotesize
\renewcommand{\arraystretch}{1.15}
\begin{tabular}{@{}rlll@{}}
\toprule
\textbf{Step} & \textbf{Outer $t_L$} & \textbf{Inner absorbed} & \textbf{\code{cand} after}\\
\midrule
1 & 2  & 1               & 1\\
2 & 7  & 5, then 6       & 6\\
3 & 11 & 10              & 10\\
4 & 12 & --- (next is 14, fails $t_R \le 12$) & 10\\
\bottomrule
\end{tabular}
\caption{Trace of Figure~\ref{fig:merge}. Total inner advances: 5. A LATERAL plan would
perform 4 independent index descents instead.}
\end{table}

\paragraph{The strict variant.} Running the same data with $<$ in place of $\le$ changes
nothing in Figure~\ref{fig:merge}, because no outer timestamp coincides with an inner one.
That is precisely why test data must be built to contain exact coincidences: give the inner
side a row at $t = 7$ and the two variants diverge at outer row $7$ --- non-strict absorbs it
and answers $7$, strict stops at $6$. A single comparison operator in loop~(D) separates the
two, and no amount of testing on non-coinciding data will tell you which one you compiled.

\paragraph{What the trace shows about cost.} Four outer rows consumed five inner advances,
and the cursor moved monotonically. A LATERAL plan over the same data performs four
independent index descents --- each starting from the root, each a fresh random access ---
and learns nothing from one to the next. The gap widens with $n$: the merge pass does
$O(n + m)$ \emph{sequential} steps in total, while the LATERAL plan does $n$ \emph{random}
descents of $O(\log m)$ each. Chapter~\ref{ch:planner} explains how the planner is taught to
reach for the former.


% ============================================================
\chapter{Making the Planner Choose It}\label{ch:planner}
% ============================================================

A correct executor node is useless if no \code{Path} leads to it. This chapter covers the
three things the planner must learn: the required sort order, the legality constraints, and
the cost.

\section{Pathkeys: the planner's model of order}
\label{sec:pathkeys}

The merge algorithm's precondition is that both inputs arrive sorted ascending on
$(k_1, \dots, k_j, t)$ --- \emph{equality keys first, match column last}. In planner terms
this is a list of \code{PathKey}s, and the mechanism for obtaining it is the same one merge
join uses:

\begin{enumerate}[leftmargin=*]
\item Build the required \code{PathKey} list for each side with the
      \code{make\_pathkeys\_for\_sortclauses} family, using the \code{btree} opfamily
      established during Phase~1 validation.
\item For each candidate child \code{Path}, compare its \code{pathkeys} against the
      requirement with \code{pathkeys\_contained\_in()}. A path whose existing order has the
      requirement as a \emph{prefix} needs no sort.
\item Otherwise wrap the child in \code{create\_sort\_path()}, paying the sort cost.
\end{enumerate}

\begin{keyidea}
This step is where the performance claim is cashed. An index on
\code{quotes (symbol, ts)} produces an \code{Index Scan} path whose pathkeys already match,
so \emph{both} sorts disappear and the join degenerates to a pure $O(n+m)$ streaming pass
over two index scans. Without the pathkey machinery we would sort unconditionally and give
back much of the advantage.
\end{keyidea}

Note the ordering requirement is on $(k, t)$ with $k$ \emph{leading}: the algorithm needs
key groups to be contiguous (remark (B) of \S\ref{sec:candidate}). An index on
\code{(ts, symbol)} does not help, which is a useful thing to be able to explain to a user
staring at an unexpected \code{Sort} node.

\section{Path generation}

\code{add\_paths\_to\_joinrel()} in \file{optimizer/path/joinpath.c} is where each join
strategy proposes its paths. For \code{JOIN\_ASOF} the \asof{} path is the \emph{only}
candidate:

\begin{itemize}[nosep]
\item \textbf{no hash join} --- hashing implements equality, and the match condition is an
      inequality;
\item \textbf{no nested loop} --- it would have to re-search per outer row, which is the
      very thing we are replacing (and the Phase~1 rewrite already covers that plan shape);
\item \textbf{merge-style only}, with sorts added as needed.
\end{itemize}

The path itself is created by a new \code{create\_asofjoin\_path()} in
\file{optimizer/util/pathnode.c}, modelled closely on \code{create\_mergejoin\_path()}.

\section{Why ASOF constrains join order}

Property (P2) --- non-commutativity --- has teeth in the optimiser. PostgreSQL's planner
explores join orders freely for inner joins, because
$A \bowtie B \equiv B \bowtie A$ and
$(A \bowtie B) \bowtie C \equiv A \bowtie (B \bowtie C)$. Outer joins already break those
identities, and \file{optimizer/path/joinrels.c} contains the machinery that restricts
them: \code{join\_is\_legal()} decides whether a given pair may be joined at all, and
\code{have\_join\_order\_restriction()} tells the planner not to give up on a join order it
would otherwise prune.

\asof{} must join the same restrictions:

\begin{itemize}[nosep]
\item \textbf{Not commutable.} $L$ must remain the outer side. The identity of ``which side
      is probed'' is part of the operator's meaning, not an implementation detail.
\item \textbf{Not freely reassociable.} Pushing another join inside an \asof{} join can
      change which rows exist on the inner side and therefore which row is nearest ---
      the same hazard that makes outer-join reassociation conditional.
\item \textbf{Qual placement.} A filter on the inner relation changes the candidate set,
      hence the answer. It must not be pushed below the \asof{} join unless it was written
      there, exactly as for the nullable side of an outer join.
\end{itemize}

The practical approach is to find every \code{switch (jointype)} in the optimiser and
handle \code{JOIN\_ASOF} deliberately. Building with \code{-Wswitch} turns this from an
exercise in grep into a compiler-enforced checklist --- a worthwhile technique in general
when adding a variant to a widely-switched enum.

\section{The cost model}

PostgreSQL costs a path in two parts: \code{initial\_cost\_xxx} (cheap, used to prune) and
\code{final\_cost\_xxx} (exact, for survivors). For \asof{}:

\begin{align}
\text{startup} &= C_{\text{sort}}(L) + C_{\text{sort}}(R) \quad\text{(either term zero if presorted)}\\
\text{total}   &= \text{startup} + (n + m)\,c_{\text{tuple}} + n\,c_{\text{match}} + n\,c_{\text{proj}}
\end{align}

where $c_{\text{match}}$ is the cost of evaluating the match comparison and the key
equality, and $c_{\text{proj}}$ the projection cost. Three observations:

\begin{enumerate}[leftmargin=*]
\item \textbf{The sorts dominate the startup cost}, which means an \asof{} join is a poor
      choice under a \code{LIMIT} unless the inputs are presorted --- the planner needs the
      startup/total split to make that trade-off correctly. This is exactly the behaviour of
      merge join, and the reason PostgreSQL models startup cost at all.
\item \textbf{There is no rescan term.} Merge join's cost function includes an allowance for
      re-reading inner tuples after a restore; \asof{} has none
      (\S\ref{sec:candidate}), so the linear term is clean. Getting this right matters:
      over-costing the join makes the planner prefer a nested loop and the whole exercise is
      wasted.
\item \textbf{Output cardinality is known exactly} for the left variant: $n$ rows,
      by (P1). This is a rare luxury --- most join estimates are the product of a selectivity
      guess --- and it makes everything above an \asof{} join easier to plan. For the inner
      variant, estimate $n \cdot p$ where $p$ is the fraction of left rows having a
      predecessor, which can be approximated from the two relations' histograms on $t$
      (specifically, the fraction of $L$'s distribution lying above $\min t_R$ per key).
\end{enumerate}

\section{The GUC}

\code{enable\_asofjoin} (default \code{on}) is added to
\file{utils/misc/guc\_tables.c} and \file{postgresql.conf.sample}. Switching it
\code{off} reverts to the Phase~1 rewrite. This is not merely a safety valve: it is the
instrument that makes the Chapter~\ref{ch:measuring} benchmark and differential test
possible at all, since it lets the same query text run through both implementations in the
same session.


% ============================================================
\chapter{Correctness Corners}
% ============================================================

Most of the defect risk in a feature like this sits in a handful of corners. Each gets a
rule and a regression test.

\paragraph{NULL in the match column.} A NULL $t$ on either side makes the comparison
\textsc{unknown}, so the row is not a candidate. In Phase~1 this falls out of the subquery's
\code{WHERE}. In Phase~2 it must be explicit: the merge comparison has to treat NULL as
non-matching rather than as an ordering extreme. Note the trap --- in a \emph{sort}, NULLs
are ordered (nulls last by default for ascending), so they will arrive at the end of the
stream and must be skipped rather than absorbed as candidates.

\paragraph{NULL in a partition key.} Treated as never-equal, consistent with \code{=}. Both
phases get this right for free if the key comparison uses the ordinary equality operator;
the trap is using a \emph{sort} comparison instead, under which two NULLs compare equal.
This is the one place where ``sorted together'' and ``equal'' genuinely diverge, and it is
worth a dedicated test.

\paragraph{Ties on $t_R$.} Implementation-defined, documented, and tested only to the extent
that both phases agree. See the remark in Chapter~2.

\paragraph{Empty inputs.} Empty $R$: all left rows NULL-extended (left) or no rows (inner).
Empty $L$: no rows, and --- worth asserting --- the inner side should ideally not be read at
all. Both are one-line tests and both have caught real bugs in similar operators.

\paragraph{Strict versus non-strict operators.} The only difference is whether equal
timestamps match. In Phase~2 this is a single comparison operator in loop~(D); getting it
backwards produces an off-by-one-row error that is invisible unless the test data contains
exactly-equal timestamps. The test data must therefore contain them deliberately.

\paragraph{Rescan.} \code{ExecReScanAsofJoin} must reset the candidate slot, the key, and
both children. Without the candidate reset, an \asof{} join on the inner side of a nested
loop returns the \emph{previous} outer iteration's match on its first row --- a subtle,
data-dependent wrong answer. Test it under a correlated subquery.

\paragraph{Memory.} The candidate slot must live in a context that is reset per key group,
or at least not leaked per tuple. A $10^7$-row query is the right smoke test: memory should
be flat.

\paragraph{Parallelism.} The node is marked parallel-unsafe initially. A single ordered pass
over a shared inner stream does not partition trivially; the natural parallelisation is
per-key-range, which requires a partitioned sort and is out of scope. Mark it, assert it,
and note the follow-up.


% ============================================================
\chapter{Measuring It}\label{ch:measuring}
% ============================================================

\section{Reading the plans}

The three plan shapes to recognise:

\begin{lstlisting}[style=plain,caption={LATERAL workaround --- per-row index descent.},captionpos=b]
Nested Loop Left Join
  ->  Seq Scan on trades t
  ->  Limit
        ->  Index Scan Backward using quotes_symbol_ts_idx on quotes q
              Index Cond: ((symbol = t.symbol) AND (ts <= t.ts))
\end{lstlisting}

\begin{lstlisting}[style=plain,caption={Native, inputs not presorted --- two sorts, one pass.},captionpos=b]
Asof Left Join
  Match Condition: (t.ts >= q.ts)
  Asof Keys: (t.symbol = q.symbol)
  ->  Sort
        Sort Key: t.symbol, t.ts
        ->  Seq Scan on trades t
  ->  Sort
        Sort Key: q.symbol, q.ts
        ->  Seq Scan on quotes q
\end{lstlisting}

\begin{lstlisting}[style=plain,caption={Native, indexes supply the order --- the goal.},captionpos=b]
Asof Left Join
  Match Condition: (t.ts >= q.ts)
  Asof Keys: (t.symbol = q.symbol)
  ->  Index Scan using trades_symbol_ts_idx on trades t
  ->  Index Scan using quotes_symbol_ts_idx on quotes q
\end{lstlisting}

The presence or absence of the \code{Sort} nodes is the single most informative thing in
these plans, and it is why \code{EXPLAIN} output is a deliverable in its own right.

\section{Benchmark design}

\paragraph{What to vary.} Row counts $10^5 / 10^6 / 10^7$; number of distinct keys
($1$, $50$, $10^4$ --- key cardinality changes how much of the inner stream each outer row
skips); index present or absent; left versus inner; match density (fraction of left rows
with a predecessor).

\paragraph{What to compare.} Three implementations on identical data: hand-written LATERAL,
Phase~1 rewrite (expected to be \emph{identical} --- any gap is a rewrite bug, not a
finding), and Phase~2 native with \code{enable\_asofjoin} on.

\paragraph{Pitfalls.}
\begin{itemize}[nosep]
\item \textbf{Cold cache.} First run reads from disk; subsequent runs do not. Warm up, or
      report both, but never mix them in one table.
\item \textbf{Planner noise.} Run \code{ANALYZE} before timing, and check with
      \code{EXPLAIN} that the plan is the one you think you are measuring. A benchmark that
      silently measured a different plan is the most common way to get a wrong answer here.
\item \textbf{\code{EXPLAIN ANALYZE} overhead.} Per-tuple instrumentation can add tens of
      percent, and it falls more heavily on plans with more node invocations --- i.e.\
      unfairly on the LATERAL plan. Time with \code{\textbackslash timing} for the headline
      numbers; use \code{EXPLAIN ANALYZE} for diagnosis.
\item \textbf{\code{work\_mem}.} The native plan's sorts spill to disk below a threshold and
      the measurement falls off a cliff. Report the setting.
\item \textbf{Index-only scans.} A covering index can make the LATERAL plan much faster than
      expected. Note whether the index covers the projected columns.
\end{itemize}

\section{Differential testing}

The strongest correctness tool the project has is that Phase~1 and Phase~2 compute the same
function by entirely different means. The harness runs the whole query corpus twice ---
\code{enable\_asofjoin = off}, then \code{on} --- and diffs. Any difference is a bug in one
of them.

\begin{lstlisting}
-- equivalence check, per query
SET enable_asofjoin = off;  CREATE TEMP TABLE a AS <query>;
SET enable_asofjoin = on;   CREATE TEMP TABLE b AS <query>;
SELECT count(*) FROM ((TABLE a EXCEPT ALL TABLE b)
                      UNION ALL (TABLE b EXCEPT ALL TABLE a)) d;  -- must be 0
\end{lstlisting}

Extend this with randomised inputs: generated tables with random key cardinality, random
NULL density, random tie frequency, and a randomly chosen operator. Property-based testing
finds the corners of Chapter~7 far more reliably than hand-written cases, because it does
not share the implementer's blind spots.


% ============================================================
\appendix
\chapter{Source File Map}
% ============================================================

Every PostgreSQL file the project touches, keyed to the work plan's issue IDs
(\file{plan.md}).

\begin{table}[h]
\centering
\footnotesize
\renewcommand{\arraystretch}{1.2}
\begin{tabular}{@{}p{6.6cm}p{5.3cm}p{1.5cm}@{}}
\toprule
\textbf{File} & \textbf{Role} & \textbf{Issues}\\
\midrule
\multicolumn{3}{@{}l}{\textit{Parser}}\\
\file{include/parser/kwlist.h} & keyword registration & P1-A1\\
\file{backend/parser/gram.y} & grammar productions & P1-A1/A2\\
\file{include/nodes/primnodes.h} & \code{JoinExpr} fields & P1-A3\\
\file{backend/parser/parse\_clause.c} & \code{transformFromClauseItem} & P1-B1\\
\file{backend/parser/parse\_asof.c} \textit{(new)} & the LATERAL rewrite & P1-B2\\
\addlinespace
\multicolumn{3}{@{}l}{\textit{Nodes \& deparse}}\\
\file{backend/nodes/gen\_node\_support.pl} & copy/equal/out/read generation & P1-A3, P2-C2\\
\file{backend/utils/adt/ruleutils.c} & \code{pg\_get\_viewdef} deparse & P1-A3, P2-C2\\
\addlinespace
\multicolumn{3}{@{}l}{\textit{Planner}}\\
\file{include/nodes/nodes.h} & \code{JOIN\_ASOF} in \code{JoinType} & P2-B1\\
\file{backend/optimizer/path/joinrels.c} & join legality, order restrictions & P2-B1\\
\file{backend/optimizer/path/joinpath.c} & path generation & P2-B2\\
\file{backend/optimizer/util/pathnode.c} & \code{create\_asofjoin\_path} & P2-B2\\
\file{backend/optimizer/path/costsize.c} & cost functions & P2-B3\\
\file{backend/utils/misc/guc\_tables.c} & \code{enable\_asofjoin} & P2-B3\\
\file{backend/optimizer/plan/createplan.c} & \code{create\_asofjoin\_plan} & P2-A1\\
\file{include/nodes/plannodes.h} & \code{AsofJoin} plan node & P2-A1\\
\addlinespace
\multicolumn{3}{@{}l}{\textit{Executor}}\\
\file{backend/executor/nodeAsofjoin.c} \textit{(new)} & the merge algorithm & P2-A2\\
\file{include/nodes/execnodes.h} & \code{AsofJoinState} & P2-A2\\
\file{backend/executor/execProcnode.c} & node dispatch & P2-A1\\
\file{backend/executor/execAmi.c} & \code{ExecReScan} & P2-A3\\
\addlinespace
\multicolumn{3}{@{}l}{\textit{Surface \& tests}}\\
\file{backend/commands/explain.c} & plan output & P2-C1\\
\file{test/regress/sql/asof\_join.sql} & regression suite & P1-C2\\
\file{doc/src/sgml/ref/select.sgml} & user documentation & P1-C3\\
\bottomrule
\end{tabular}
\end{table}

\chapter{Glossary}

\begin{description}[leftmargin=*,style=nextline,font=\normalfont\bfseries\ttfamily]
\item[Path] One candidate way of producing a relation, with cost and sort order. The
planner enumerates paths and keeps the cheapest.
\item[PathKey] Canonical representation of a sort order, used to decide whether a sort can
be skipped.
\item[RelOptInfo] Planner's record of a relation --- base or join --- and its paths.
\item[TupleTableSlot] Holder for one tuple in the executor. Reused by the producing node,
hence the copy requirement in \S\ref{sec:candidate}.
\item[LATERAL] \code{FROM}-clause modifier permitting a subquery to reference earlier
\code{FROM} items; evaluated per outer row.
\item[Desugaring] Defining a new construct by translation into existing ones. Phase~1's
technique.
\item[Volcano model] Iterator-based execution: every node exposes \code{next()}, pulling
tuples from children on demand.
\item[Mark/restore] Merge join's mechanism for rewinding the inner side over duplicate
keys. Not needed by \asof{} --- see \S\ref{sec:candidate}.
\item[opfamily] Catalogued group of related operators (e.g.\ \code{btree} comparison
operators for a type), which is how PostgreSQL knows a type has a usable total order.
\item[GUC] ``Grand Unified Configuration'' --- a PostgreSQL run-time setting, e.g.\
\code{enable\_asofjoin}.
\end{description}

\end{document}
